\vfill\pagebreak
\section{Patterns that change what they match}
\label{sec:mutability:patterns}

A pattern declares variables: the names of \token{(d, m)} or of
\token{Ok(n)} (Section~\ref{sec:collections:match}). Like the iterators of a
loop, they are constant copies of the parts of the value matched. They can also
be declared as a loop iterator can, with \token{dmut} for a part whose elements
change, and with \token{ref mut} for a reference to the part itself. The value
matched then carries \token{alias} or \token{ref}, as the argument of a call
does.

\subsection{\texttt{match alias}}

A variable of a pattern that shares elements that it may change is declared
\token{dmut}, and the value matched is given with \token{alias}.

\begin{lstlisting}[style=coloredverbatim]
let dmut team = (1, copy [12, 15, 9]);
match alias team {
  (1, dmut scores) => {
    scores[0] = 20;
  }
  _ => {}
}
println(team);
\end{lstlisting}
\vspace{-10pt}%
\begin{lstlisting}[style=bashVerb]
(1, [20, 15, 9])
\end{lstlisting}

\token{scores} is a copy of the second field of \token{team}, a length and an
address, and shares its elements as the tuples of
Section~\ref{sec:mutability:alias} did. Without \token{alias}, the pattern is
refused with \errmsg{E4045}{discarding the constant property is prohibited},
and the note under it names the field, \textit{for pattern field access 1 on
  type (mut i32, mut [mut i32])}.

\subsection{\texttt{match ref}}

A variable of a pattern declared \token{ref mut} is a reference to a part of
the value matched, and assigning it writes into that part. The value matched is
given with \token{ref}, and must be a variable, or a part of one. The next
listing puts the two bounds of an interval in order.

\begin{lstlisting}[style=coloredverbatim]
let dmut range = (9, 4);
match ref range {
  (ref mut low, ref mut high) if low > high => {
    let old = low;
    low = high;
    high = old;
  }
  _ => {}
}
println(range);
\end{lstlisting}
\vspace{-10pt}%
\begin{lstlisting}[style=bashVerb]
(4, 9)
\end{lstlisting}

\token{low} and \token{high} refer to the two fields of \token{range}, and the
guard reads them as it would read copies. \token{range} is declared
\token{dmut}: its fields are a level of their own, as the elements of an array
are, and \token{let mut} would let the program give \token{range} a new tuple,
but not change one field.

The same works on an option, whose value can then be changed in place, and on
a slice, whose pattern can take references to its first elements.

\begin{lstlisting}[style=coloredverbatim]
let dmut best: i32? = none;
for score in copy [12, 17, 9] {
  if let Ok(ref mut b) = ref best {
    if score > b {
      b = score;
    }
  } else {
    best = score?;
  }
}
println(best);

let dmut word = copy "hello";
match ref word {
  [ref mut first, _...] => {
    first = 'H';
  }
  _ => {}
}
println(word);
\end{lstlisting}
\vspace{-10pt}%
\begin{lstlisting}[style=bashVerb]
Ok(17)
Hello
\end{lstlisting}

A \token{let} declaration is a pattern too, and takes the same keywords. The
next declaration makes two references to the fields of \token{range}.

\begin{lstlisting}[style=coloredverbatim]
let dmut range = (9, 4);
let (ref mut low, ref mut high) = ref range;
low = 0;
println(range, " ", high);
\end{lstlisting}
\vspace{-10pt}%
\begin{lstlisting}[style=bashVerb]
(0, 4) 4
\end{lstlisting}

\subsection{Which keyword}

A pattern can mix the three kinds of variables: constant copies, \token{dmut}
copies that share elements, and \token{ref mut} references. The keyword in front
of the value matched depends on the variables of all the arms.

\begin{itemize}
\item \token{ref}, if a variable is declared \token{ref mut}. It also covers the
  \token{dmut} variables of the same pattern, as in
  \token{[ref mut first, dmut rest...]}, which takes a reference to the first
  element and shares the others.
\item \token{alias}, if a variable is declared \token{dmut} and shares elements
  that it may change, and none is a reference that may change.
\item Neither, if every variable is a constant. \token{match ref range} with the
  pattern \token{(a, b)} is refused with \errmsg{E4252}{referencing the value
    has no effect}, the note under it saying \textit{there is no pattern that
    requires the value to be passed by mutable reference}, and
  \token{match alias} is refused in the same case with the
  \errmsg{E4247}{aliasing the value \errhole{} to create a constant borrowing
    is prohibited} of Section~\ref{sec:mutability:alias}.
\end{itemize}

A variable declared \token{ref mut} when the value is matched without
\token{ref} is refused with \errmsg{E4089}{ref of type mut i32 must be mutable,
  but is const}: without the keyword, the parts of the value are constant for
the pattern.
